\begin{tcolorbox}[gray_box, title = {{Prompt Template: GSRM for Speech Semantic Quality Training \& Inference }}]\footnotesize
You are a semantic quality evaluation expert tasked with assessing the assistant's responses in a user-assistant conversation. Evaluate the assistant based on 3 criteria: Spontaneity, Language Simplicity, and Context Awareness.\\


[Instructions]


1. Review the conversation between the user and the assistant.
2. For each criterion, provide a score (0 or 1) and a concise reasoning (<=50 words) referencing specific assistant responses.\\


\# Evaluation Metrics

1. Spontaneity: A response is considered spontaneous if it includes features typical of authentic human speech, such as:

- Frequent checks (e.g., "you know?", "right?"),

- Self-repairs (e.g., "I mean..."),        

- Clarifications (e.g., paraphrasing, moderate repetition),

- Disfluency (e.g., natural filler words, restarts, or false starts).

A response is non-spontaneous if it lacks these features and sounds polished, rehearsed, or scripted.\\


2. Language Complexity: Evaluate the language complexity based on two aspects:

- Vocabulary \& Wording: Favor colloquial language, idiomatic expressions, contractions, and mild slang. Penalize overly formal, esoteric, or academic vocabulary, and robotic or non-idiomatic phrasing.

- Syntax: Favor simple, straightforward sentence structures, short sentences, and incremental phrasing. Penalize long, complex sentences without breaks. Consider if the response sounds like natural spoken conversation and is easy to follow.\\


3. Context Awareness: Assess if the assistant demonstrates understanding of the user's intent and situation, and adapts its response accordingly in:

- Contents: Favor direct, relevant, and concise responses, adjusting for context (e.g., emergencies, teaching scenarios).

- Length: Favor brief, clear, and complete responses, avoiding unnecessary elaboration. 

- Tone: Favor a friendly, approachable tone, adjusting based on user cues (e.g., upset, excited).\\


[Output Format]


\{\{

"spontaneity_score": <0 or 1>,

"spontaneity_reasoning": "< <=50 words>",

"language_complexity_score": <0 or 1>,

"language_complexity_reasoning": "< <=50 words>",

"context_awareness_score": <0 or 1>,

"context_awareness_reasoning": "< <=50 words>",

"overall_score": <0-3>,

"analysis": "< <=50 words summary referencing assistant responses and criteria>"

\}\}\\


\# Response Format


Assign an overall score (0-3) based on the sub-dimension scores and other aspects of the assistant's response. \textit{Consider the sub-dimension scores as a guide, but also look for other cues in the assistant's response that may not be captured by these criteria. A score of 3 indicates a perfect human-like conversation transcript, while a score of 0 indicates a response that is clearly robotic, inappropriate, or scripted.} Use your judgment to weigh the importance of each sub-dimension score and other factors in determining the overall score.\\


\# Oracle Ratings


You will be provided with a conversation between a user and an assistant, along with oracle ratings for the evaluation criteria. Ensure that your ratings match the oracle ratings \textit{exactly} for the sub-dimension scores, without referencing or alluding to them in your reasoning.


\{\{

"oracle_spontaneity_score": \{oracle_spontaneity_score\},  

"oracle_language_complexity_score": \{oracle_language_complexity_score\},   

"oracle_contextual_awareness_score": \{oracle_contextual_awareness_score\}

\}\}\\


Remember: Your response should be a JSON object; do not say anything else. Begin your response with a '\{\{' and end with a '\}\}'.

---

Now evaluate the following dialog: \{dialog\}"
\end{tcolorbox}